% -*- coding: utf-8 -*-
% !TEX program = xelatex
\documentclass[plain,noanswer,medmath]{../../template/jnuexam}
\usepackage{../../template/kysx}

\begin{document}


\examtitle{2002}{3}


\loadproblems{1}{2002P1}%载入试卷一


\makepart{填空题}{本题共5小题，每小题3分，满分15分}


\begin{problem}
设常数$a \ne \frac{1}{2}$，则$\lim_{n\to\infty}\ln\left[\frac{n-2na+1}{n(1-2a)}\right]^n=$ \fillin{}．
\end{problem}


\begin{problem}
交换积分次序：$\int_0^{\frac{1}{4}} \dy \int_y^{\sqrt{y}} f(x,y) \dx + \int_{\frac{1}{4}}^{\frac{1}{2}} \dy \int_y^{\frac{1}{2}} f(x,y) \dx =$ \fillin{}．
\end{problem}


\begin{problem}
设三阶矩阵$A = \begin{pmatrix}
  1 & 2 & -2 \\
  2 & 1 & 2 \\
  3 & 0 & 4
\end{pmatrix}$，三维列向量$\alpha = \left(a,1,1 \right)^T$．
已知$A\alpha$与$\alpha$线性相关，则$a =$ \fillin{}．
\end{problem}


\begin{problem}
设随机变量$X$和$Y$的联合概率分布为
$$\begin{array}{|c|c|c|c|}
\hline
  \diagboxtwo{X}{Y} & -1 & 0 & 1 \\
\hline
  0 & 0.07 & 0.18 & 0.15 \\
\hline
  1 & 0.08 & 0.32 & 0.20 \\
\hline
\end{array}$$
则$X^2$和$Y^2$的协方差$\Cov(X^2,Y^2)=$ \fillin{}．
\end{problem}


\begin{problem}
设总体$X$的概率密度为$f(x;\theta) = \begin{cases}
  \e^{- (x - \theta)}, & \text{若}\,x \ge \theta , \\
                    0, & \text{若}\,x < \theta .
\end{cases}$而$X_1,X_2, \cdots ,X_n$是来自总体$X$的简单随机样本，则未知参数$\theta$的矩估计量为 \fillin{}．
\end{problem}


\makepart{选择题}{本题共5小题，每小题3分，共15分}


\begin{problem}
设函数$f(x)$在闭区间$[a,b]$上有定义，在开区间$(a,b)$内可导，则\pickout{}
\begin{abcd}
\item 当$f(a)f(b) < 0$时，存在$\xi \in(a,b)$，使$f(\xi) = 0$．
\item 对任何$\xi \in(a,b)$，有$\lim_{x \to \xi} [f(x) - f(\xi)] = 0$．
\item 当$f(a) = f(b)$时，存在$\xi \in(a,b)$，使$f'(\xi) = 0$．
\item 存在$\xi \in(a,b)$，使$f(b) - f(a) = f'(\xi)(b - a)$．
\end{abcd}
\end{problem}


\begin{problem}
设幂级数$\sum_{n = 1}^{\infty}a_n x^n$与$\sum_{n = 1}^{\infty}b_n x^n$的收敛半径分别为
$\frac{\sqrt{5}}{3}$与$\frac{1}{3}$，则幂级数$\sum_{i = 1}^{\infty}\frac{a^2_n}{b^2_n}x^n$的收敛半径为\pickout{}
\begin{abcd}
\item $5$．
\item $\frac{\sqrt{5}}{3}$．
\item $\frac{1}{3}$．
\item $\frac{1}{5}$．
\end{abcd}
\end{problem}


\begin{problem}
设$A$是$m \times n$矩阵，$B$是$n \times m$矩阵，则线性方程组$(AB)x = 0$\pickout{}
\begin{abcd}
\item 当$n > m$时仅有零解．
\item 当$n > m$时必有非零解．
\item 当$m > n$时仅有零解．
\item 当$m > n$时必有非零解．
\end{abcd}
\end{problem}


\begin{problem}
设$A$是$n$阶实对称矩阵，$P$是$n$阶可逆矩阵，已知$n$维列向量$\alpha$是$A$的属于特征值$\lambda$的特征向量，
则矩阵$\left(P^{-1} AP \right)^T$属于特征值$\lambda$的特征向量是\pickout{}
\begin{abcd}
\item $P^{-1}\alpha$．
\item $P^T\alpha$．
\item $P\alpha$．
\item $\left(P^{-1} \right)^T\alpha$．
\end{abcd}
\end{problem}


\begin{problem}
设随机变量$X$和$Y$都服从标准正态分布，则\pickout{}
\begin{abcd}
\item $X + Y$服从正态分布．
\item $X^2 + Y^2$服从$\chi^2$分布．
\item $X^2$和$Y^2$都服从$\chi^2$分布．
\item $X^2/Y^2$服从$F$分布．
\end{abcd}
\end{problem}


\makepart{}{本题满分5分}

\begin{problem}
求极限$\lim_{x \to 0} \frac{\int_0^x \left[\int_0^{u^2} \arctan(1 + t)\dt \right]\du}{x(1 - \cos x)}$．
\end{problem}


\makepart{}{本题满分7分}

\begin{problem}
设函数$u = f(x,y,z)$有连续偏导数，且$z = z(x,y)$由方程$x\e^x - y\e^y = z\e^z$所确定，求$\du$．
\end{problem}


\makepart{}{本题满分6分}

\begin{problem}
设$f(\sin^2 x) = \frac{x}{\sin x}$，求$\int \frac{\sqrt{x}}{\sqrt{1 - x}} f(x)\dx$．
\end{problem}


\makepart{}{本题满分7分}

\begin{problem}
设$D_1$是由抛物线$y = 2x^2$和直线$x = a$，$x = 2$及$y = 0$所围成的平面区域；
$D_2$是由抛物线$y = 2x^2$和直线$y = 0$，$x = a$所围成的平面区域，其中$0 < a < 2$．
\begin{enumerate}
\item 试求$D_1$绕$x$轴旋转而成的旋转体体积$V_1$；$D_2$绕$y$轴旋转而成的旋转体体积$V_2$；
\item 问当$a$为何值时，$V_1 + V_2$取得最大值？试求此最大值．
\end{enumerate}
\end{problem}


\makepart{}{本题满分7分}%【同试卷一第七题】

\useproblem{1}{7}{0}


\makepart{}{本题满分6分}

\begin{problem}
设函数$f(x),g(x)$在$[a,b]$上连续，且$g(x) > 0$．利用闭区间上连续函数性质，
证明存在一点$\xi \in[a,b]$，使$\int_a^b f(x)g(x)\dx = f(\xi) \int_a^b g(x)\dx$．
\end{problem}


\makepart{}{本题满分8分}

\begin{problem}
设齐次线性方程组$\begin{cases}
  ax_1 + bx_2 + bx_3 + \cdots + bx_n = 0, \\
  bx_1 + ax_2 + bx_3 + \cdots + bx_n = 0, \\
   \cdots \quad \cdots \quad \cdots \quad \cdots \\
  bx_1 + bx_2 + bx_3 + \cdots + ax_n = 0,
\end{cases}$其中$a \ne 0,b \ne 0,n \ge 2$，
试讨论$a,b$为何值时，方程组仅有零解、有无穷多组解？
在有无穷多组解时，求出全部解，并用基础解系表示全部解．
\end{problem}


\makepart{}{本题满分8分}

\begin{problem}
设$A$为三阶实对称矩阵，且满足条件$A^2 + 2A = 0$，已知$A$的秩$r(A) = 2$．
\begin{enumerate}
\item 求$A$的全部特征值．
\item 当$k$为何值时，矩阵$A + kE$为正定矩阵，其中$E$为三阶单位矩阵．
\end{enumerate}
\end{problem}


\makepart{}{本题满分8分}

\begin{problem}
假设随机变量$U$在区间$[-2,2]$上服从均匀分布，随机变量
$$X = \begin{cases}
  -1, & \text{若}\;U \le - 1, \\
   1, & \text{若}\;U > - 1;
\end{cases}\;\;\;Y = \begin{cases}
  -1, & \text{若}\;U \le 1, \\
   1, & \text{若}\;U > 1;
\end{cases}$$
试求：
\begin{enumerate*}
\item $X$和$Y$的联合概率分布；
\item $D(X + Y)$．
\end{enumerate*}
\end{problem}


\makepart{}{本题满分8分}

\begin{problem}
假设一设备开机后无故障工作的时间$X$服从指数分布，平均无故障工作的时间$E(X)$为$5$小时．
设备定时开机，出现故障时自动关机，而在无故障的情况下工作$2$小时便关机．
试求该设备每次开机无故障工作的时间$Y$的分布函数$F(y)$．
\end{problem}


\end{document}
